\subsection{Second-order consistency on residual boards}\label{sec:B4T05}
\begin{teorema}[Macroscopic consistency]\label{thm:macro}
Fix $\delta_0,C,\eta,\gamma>0$. Let $B$ be a realizable ordinary residual board with $k\ge\delta_0n$ empty
rows and as many empty columns. Suppose its Bethe point $p$ is stationary for $H_B$ and lies in the
interior, with all RC masses equal to one and
\[
 p_e\le C/k,\qquad h_D\ge\eta,\qquad
 \lambda_2(\mathscr L_p)\ge\gamma.
\]
The RC graph is connected. The predictor uses the RC cofactor of the canonical determinant, the local
factor, and $\tfrac12\log(2\pi)$ with the conventions of Sections~\ref{sec:B3T01} and~\ref{sec:B4T01}. Then
there exist $k_0$ and $C_*$ depending only on these four constants such that, for $k\ge k_0$, every legal
child has a unique interior stationary point on the small branch and
\[
 \boxed{\left|L(B+c)-L(B)-\log[p_c(1+\delta_c)]\right|
 \le C_*/k^2\quad\hbox{for every child }c.}
\]
Here $\delta_c=(R\nabla\ell)_c/p_c=O(k^{-1})$. No unfavorable children are excluded, and no per-square lower
bound is required.
\end{teorema}
\begin{proof}
The existence and the bounds for the child stationary point are Lemma~\ref{lem:hijo}.
Section~\ref{sec:B4T09}
combines the action and the energy of Sections~\ref{sec:B4T04} and~\ref{sec:B4T06} with
Lemmas~\ref{lem:prefactor} and~\ref{lem:modocero}, and its final result is the bound on $r_c$, which by
Lemma~\ref{lem:escalar} is equivalent to $\mathcal C_c=O(k^{-2})$. The table in that section collects the
terms of order $1/k$ and their cancellation.
\end{proof}

\begin{corolario}\label{cor:rho}
Under the hypotheses of Theorem~\ref{thm:macro}, for each empty row $r$ we have $|\log\rho(B)|\le C_*/k^2$,
where $\rho(B)=\sum_{c\in r}e^{L(B+c)-L(B)}$.
\end{corolario}
\begin{proof}
The differential identity of Section~\ref{sec:B3T01} gives a sum of one in each row.
Positivity follows from $p_c>0$ and from $|\delta_c|=O(k^{-1})$ for large $k$ (Section~\ref{sec:B3T06}).
Thus,
\[
 w_c=p_c(1+\delta_c)>0,\qquad\sum_{c\in r}w_c=1,
 \quad \rho(B)=\sum_{c\in r}e^{L(B+c)-L(B)}=\mathbb E_w e^r.
\]
Therefore $\min r\le\log\rho\le\max r$, which implies
\[
 \boxed{|\log\rho(B)|\le C_*/k^2.}
\]
\end{proof}
If a trajectory of parents stays in the same quantitative class, with the same constants and
$k\ge\delta_0n$, then
\[
 \sum_{k\ge\delta_0n}|\log\rho(B_k)|\le C_*\sum_{k\ge\delta_0n}k^{-2}
 =O_{\delta_0,C,\eta,\gamma}(n^{-1}).
\]
Membership of the entire trajectory in the class is a separate hypothesis. We do not iterate the parameter
loss from one-step inheritance without controlling it.

\begin{corolario}[The full board]\label{cor:completo}
There exist $n_1$ and $C_*$ such that, for every $n\ge n_1$ and every row $r$ of the full $n\times n$
board $B$, we have $|\log\rho(B)|\le C_*/n^2$.
\end{corolario}
\begin{proof}
For $n\ge n_0$, (P2.7) gives $c/n\le\pi_e\le M/n$ and $h_d\ge\eta$, and the canonical maximizer is the
interior stationary point of $H_B$ with RC masses one. As a quadratic form, the Laplacian $\mathscr L_\pi$
is at least $c/n$ times that of the complete bipartite graph $K_{n,n}$, because the difference is the
Laplacian of a graph with nonnegative weights. The second eigenvalue of the Laplacian of $K_{n,n}$ is
$n$, so $\lambda_2(\mathscr L_\pi)\ge c$. Theorem~\ref{thm:macro} and Corollary~\ref{cor:rho} apply with
$k=n$, $\delta_0=1$, $C=M$ and $\gamma=c$.
\end{proof}
The class of Theorem~\ref{thm:macro} is therefore not empty, and on the full board the predictor is
$Q_{\rm pred}(n)$ (Lemma~\ref{lem:puente}).

\paragraph{Scope.}
The theorem is deterministic within the stated class. We do not study here whether typical trajectories
remain in it.
